% mooc_SOTAai_B.tex
\begin{center}
  \Large
  \color{colorA}
  \textbf{Módulo 2:\\ Generación de listado inicial de palabras clave}
\end{center}

\begin{table}[ht!]
  \centering
  \def\arraystretch{1.8}
  \caption{Acotación del tema y el periodo.}
  \begin{tabular}{L{0.75\textwidth}c}
    \toprule
    \bfseries Tema & \bfseries Periodo \\ \hline
    El uso de la metaheurística BRKGA en la asignación de personal
    por turnos, o jornada completa. Formulado como un problema de optimización multi-objetivo sujeto a gran cantidad de restricciones laborales y operativas.
    & 2020-2026 \\
    \bottomrule
  \end{tabular}
\end{table}

\begin{table}[ht!]
  \centering
  \def\arraystretch{1.8}
  \caption{Tabla de sinónimos.}
  \begin{tabularx}{\textwidth}{W{1.4}W{1.0}W{0.6}}
    \toprule
    \bfseries Concepto & \bfseries Variantes / Acrónimos & Notas \\ \hline
    Algoritmos metaheurísticos poblacionales & BRKGA, GA & \\
    Algoritmos genéticos de llaves aleatorias sesgadas & Biased random-key genetic algorithm / BRKGA & \\
    Timetabling & & \\
    Scheduling & & \\
    Personnel rostering &Staff rostering & \\
    Workforce scheduling metaheuristic& & \\
    Air traffic controller scheduling & ATCO scheduling & \\
    controller-to-sector assignment & ATCO assign & \\
    Multi-objective shift scheduling& & \\
    Fatigue-aware scheduling& & \\
    % & & \\
    % & & \\
    \bottomrule
  \end{tabularx}
\end{table}

% \href{run:Users/ericlucero/Documents/Academic/Courses_Seminars/MOOCS_UPM/SOTAai/M2_Guia_Operativa.pdf}{pdf}
\href{file:///Users/ericlucero/Documents/Academic/Courses_Seminars/MOOCS_UPM/SOTAai/M2_Guia_Operativa.pdf}{pdf}
\begin{table}[ht!]
  \centering
  \def\arraystretch{1.8}
  \caption{Tabla de resultados.}
  \begin{tabularx}{\textwidth}{|W{0.65}|W{0.35}|W{1.0}|W{1.0}|W{1.35}|W{1.65}|}
    \hline
    \bfseries Fuente & \bfseries Año& \bfseries DOI& \bfseries Título& \bfseries Abstract& \bfseries Cadena de búsqueda \\ \hline
    Scopus& 2026& & & & \\
    Scopus& & & & & \\
    PubMed & & & & & \\
    PubMed & & & & & \\
    ArXiV & & & & & \\
    ArXiV & & & & & \\
    CrossRef & & & & & \\
    % & & \\
    \hline
  \end{tabularx}
\end{table}